% Copy flat_fixed_dynamic_tree_cd_overlap_averages.pdf into your Overleaf
% project's figures/ directory, then use this figure block in the paper.
% The paragraph below is commentary for the main text. The caption only describes
% the visual encoding in the figure.

This sample-level overlap analysis shows that all structured memory methods
recover cross-domain samples that fail under the Flat Memory List, so the
improvement is not only visible in aggregate failure rates. Averaged over the
seven evaluated models, dynamic partitions remain strongest overall, recovering
the largest average number of Flat-failing samples while keeping shared failures
below fixed partitions and the two-level tree. Compared with fixed partitions,
dynamic partitions recover more samples where the Flat Memory List fails and
also reduce the average number of samples where both methods fail.

\begin{figure}[t]
  \centering
  \includegraphics[width=\linewidth]{figures/flat_fixed_dynamic_tree_cd_overlap_averages.pdf}
  \caption{
  Average sample-level cross-domain outcome overlap between the Flat Memory List
  and each structured memory method, averaged over the seven evaluated models.
  Each horizontal bar compares the Flat Memory List against one method. Green
  denotes samples where both methods pass; blue denotes samples where the Flat
  Memory List passes but the comparison method fails; orange denotes samples
  where the comparison method passes but the Flat Memory List fails; and red
  denotes samples where both methods fail.
  }
  \label{fig:flat-fixed-dynamic-tree-cd-overlap-averages}
\end{figure}
